% Agent 10: Book pages 381--384
% Section 2.7 Shock Waves (continued), Section 2.8 Turbulence,
% Section 2.9 Reconnection of Magnetic Field Lines (beginning)
% Equations (2.7.1a)--(2.7.10), (2.9.1)--(2.9.5)

%%%%%%%%%%%%%%%%%%%%%%%%%%%%%%%%%%%%%%%%%%%%%%%%%%%%%%%%%%%%%%%%
%%% SECTION 2.7  SHOCK WAVES  (continued from previous agent)
%%%%%%%%%%%%%%%%%%%%%%%%%%%%%%%%%%%%%%%%%%%%%%%%%%%%%%%%%%%%%%%%

% --- The relativistic Rankine--Hugoniot equations ---
% (Section heading and basic references on previous agent's pages)

\subsubsection*{The relativistic Rankine--Hugoniot equations}

Pick a particular event~$\mathscr{P}$ on a shock front.  In the neighborhood
of~$\mathscr{P}$ introduce a local Lorentz frame (``rest frame of the
shock'') in which (i)~the shock is the surface $y = z = 0$; and (ii)~on both
sides of the shock the fluid is moving in the $x$~direction, i.e.,
perpendicular to the shock front (``normal shock''; see Figure~2.7.1).
That such a local Lorentz frame exists in general one can prove quite easily
[see, e.g., \citet{Taub1948}].

Denote the ``front'' side of the shock (side \emph{from} which the fluid
moves) by a ``1'', and denote the ``back'' side (side \emph{toward} which the
fluid moves) by a ``2''; see Figure~2.7.1.  Denote the velocity of the fluid,
as measured in the rest frame of the shock, by
%
\begin{align}
  v_1 &= (dx/dt)_1 = \text{ordinary velocity on front side,}
  \tag{2.7.1a} \\
  v_2 &= (dx/dt)_2 = \text{ordinary velocity on back side,}
  \tag{2.7.1b}
\end{align}
%
\begin{align}
  u_1 &= v_1\gamma_1 = v_1/(1 - v_1^{2})^{1/2}\,,
  \tag{2.7.1c} \\
  u_2 &= v_2\gamma_2 = v_2/(1 - v_2^{2})^{1/2}\,.
  \tag{2.7.1d}
\end{align}
%
(Note that these ``4-velocities'' are scalars, not vectors.)

In the rest frame of the shock the law of baryon conservation is equivalent
to continuity of the \emph{baryon flux}:
%
\begin{equation}
  j = n_1 u_1 = n_2 u_2\,.
  \tag{2.7.2a}
\end{equation}

\begin{figure}[htbp]
\centering
\fbox{\parbox{0.9\columnwidth}{\centering [Figure 2.7.1 placeholder]\\
A shock wave viewed in the local-Lorentz rest frame of the shock front at an
event~$\mathscr{P}$.}}
\caption{A shock wave viewed in the local-Lorentz rest frame of the shock
  front at an event~$\mathscr{P}$.}
\label{fig:2.7.1}
\end{figure}

\noindent
Similarly, energy and momentum conservation are equivalent to continuity of
energy flux and momentum flux:
%
\begin{align}
  \bigl[(\varrho + p)\,v\,u\bigr]_1
    &= \bigl[(\varrho + p)\,v\,u\bigr]_2\,,
  \tag{2.7.2b} \\
  \bigl\{(\varrho + p)\,u^2 + p\bigr\}_1
    &= \bigl\{(\varrho + p)\,u^2 + p\bigr\}_2\,.
  \tag{2.7.2c}
\end{align}

These junction conditions are more easily understood by writing them in a
form analogous to the Rankine-Hugoniot equations of Newtonian theory: First
take the law of baryon conservation~(2.7.2a) and turn it into equations for
the fluid 4-velocities
%
\begin{equation}
  u_1 = jV_1\,, \qquad u_2 = jV_2\,.
  \tag{2.7.3a}
\end{equation}
%
(See \S2.5 for notation.) Then take the law of momentum conservation~(2.7.2c),
rewrite it in terms of $\mu$, $V$, $j$, and~$p$ using the law of baryon
conservation~(2.7.2a) and the relations $n = 1/V$,
$\mu = (\varrho + p)\,V$; and solve for the baryon flux to obtain
%
\begin{equation}
  j^2 = -\frac{p_2 - p_1}{\mu_2 V_2 - \mu_1 V_1}\,.
  \tag{2.7.3b}
\end{equation}

Finally, use the law of baryon conservation~(2.7.2a) to rewrite the energy
equation~(2.7.2b) in the form
%
\begin{equation}
  \mu_1\gamma_1 = \mu_2\gamma_2\,;
\end{equation}
%
divide equation~(2.7.3b) by $(\mu_1 V_1 + \mu_2 V_2)$, and combine with
$j = u_1/V_1 = u_2/V_2$; then subtract this from the square of the energy
equation $(\mu_2\gamma_2)^2 - (\mu_1\gamma_1)^2 = 0$, to obtain
%
\begin{equation}
  \mu_2^2 - \mu_1^2 = (p_2 - p_1)(\mu_1 V_1 + \mu_2 V_2)\,.
  \tag{2.7.3c}
\end{equation}

\emph{Equations~\textnormal{(2.7.3)} are Taub's~\textnormal{(1948)} junction
conditions for shock waves.}
Their Newtonian limits are the standard Rankine-Hugoniot equations:
%
\begin{align}
  v_1 &= j_0\,V_{01}\,, \qquad v_2 = j_0\,V_{02}\,,
  \tag{2.7.4a} \\[4pt]
  \hat{\jmath}_0^{\,2} &\equiv \frac{p_2 - p_1}{V_{02} - V_{01}}
  = \text{``mass flux'' of Newtonian theory}\,;
  \tag{2.7.4b} \\[4pt]
  w_2 - w_1 &= \tfrac{1}{2}(p_2 - p_1)(V_{01} + V_{02})\,.
  \tag{2.7.4c}
\end{align}

The form of the Newtonian junction conditions~(2.7.4) motivates one to use~$p$
and~$V_0$ as one's independent thermodynamic variables when analyzing
Newtonian shocks; similarly, the form of the relativistic junction
conditions~(2.7.3) motivates one to use~$p$ and~$\mu V$ as one's independent
variables for relativistic shocks.

%%% The Rankine-Hugoniot curve %%%

\subsubsection*{The Rankine-Hugoniot curve}

Consider a family of shocks, each with the same thermodynamic state on the
front face (same $\mu_1$, $V_1$, $p_1$, etc.), but with different states on the
back face (different $\mu_2$, $V_2$, $p_2$, etc.).  This family of shocks is
a one-parameter family.  Thus, if one plots all back-face states
$(\mu_2 V_2,\,p_2)$ in the $\mu V$--$p$ plane, they lie on a single
curve---the ``Rankine-Hugoniot curve''---passing through the point
$(\mu_1 V_1,\,p_1)$.  (In the relativistic case this curve is also called the
``Taub adiabat'' or ``Poisson adiabat''.)

One can also plot, in the $\mu V$--$p$ plane, the Poisson adiabat (curve of
constant entropy) passing through $(\mu_1 V_1,\,p_1)$.  These two curves
typically have the relative shapes and locations shown in Figure~2.7.2.  In
particular, from the Rankine-Hugoniot equations one can derive the following
general properties of the Rankine-Hugoniot
curve:\footnote{See \citet{Thorne1973a} for derivation.  (i)~The
  Rankine-Hugoniot curve is tangent to the Poisson adiabat, and has the same
  second derivative at point~``1''.  The derivation requires the assumption
  $[\partial^2(\mu V)/\partial p^2]_{s} > 0$; a condition which is satisfied
  by all materials of astrophysical or laboratory interest; see
  \citet{LandauLifshitz1959}.  Some, but not all, of the listed properties
  remain true without this condition; see \citet{ZeldovichRaizer1966} for
  nonrelativistic details, and \citet{Thorne1973a} for relativistic details.}

\begin{enumerate}
\item[(i)] The Rankine-Hugoniot curve is tangent to the Poisson adiabat, and
  has the same second derivative at point~``1''; i.e., for weak shocks the
  increase in entropy is third-order in the pressure jump:
  %
  \begin{equation}
    s_2 - s_1 = \left[\frac{1}{12\mu T}\,
      \frac{\partial^2(\mu V)}{\partial p^2}\bigg|_{s,\,1}\right]
      (p_2 - p_1)^3 + O\!\left[(p_2-p_1)^4\right].
    \tag{2.7.5}
  \end{equation}

\item[(ii)] As the gas passes from the front of the shock to the back, its
  entropy, pressure, and chemical potential increase
  %
  \begin{equation}
    s_2 > s_1\,,\qquad p_2 > p_1\,,\qquad \mu_2 > \mu_1\,;
    \tag{2.7.6a}
  \end{equation}
  %
  while its specific volume and the product $\mu V$ decrease
  %
  \begin{equation}
    V_2 < V_1\,,\qquad \mu_2 V_2 < \mu_1 V_1\,.
    \tag{2.7.6b}
  \end{equation}
  %
  (Hence, only the ``upper branch'' of the Rankine-Hugoniot curve,
  Figure~2.7.2, is physically relevant.)

\item[(iii)] The flow on the front side is always supersonic; on the back side
  is always subsonic
  %
  \begin{equation}
    u_1/u_{S1} > 1\,,\qquad u_2/u_{S2} < 1\,;
    \tag{2.7.7}
  \end{equation}
  %
  \[
    u_S = c_S/(1 - c_S^2)^{1/2}\,.
  \]

\item[(iv)] As one moves up the Rankine-Hugoniot curve, away from point
  ``1''---i.e., as one studies a sequence of ever ``stronger''
  shocks---the following quantities increase monotonically:
  %
  \begin{equation}
  \begin{gathered}
    \text{the baryon flux across the shock, } j\,; \\
    \text{the jump in entropy across the shock, } s_2 - s_1\,; \\
    \text{the ``relativistic'' Mach number on the front side of the shock, }
      M_1 = u_1/u_{S1}\,.
  \end{gathered}
    \tag{2.7.8}
  \end{equation}
\end{enumerate}

Notice that, once the thermodynamic state on the front face of the shock
has been specified, the shock has only one free parameter.  If one fixes the
baryon flux across the shock~$j$, or the speed on the front face~$u_1$, or
the pressure on the back face~$p_2$, or any other single parameter, then all
other properties of the shock are uniquely determined.  To calculate them, one
need merely invoke the Rankine-Hugoniot equations~(2.7.4), the laws of
thermodynamics, and the equation of state of the gas.

%%% Shocks in an ideal gas %%%

\subsubsection*{Shocks in an ideal gas}

Consider, as a special but important case, a nonrelativistic ideal gas with
constant adiabatic index $\Gamma = -(\partial\ln p/\partial\ln V_0)_s$, so
that
%
\begin{equation}
  pV_0 = kT/\bar{m}\,.
  \tag{2.7.9}
\end{equation}

By a straightforward calculation (\S85 of \citealt{LandauLifshitz1959}), one
can derive the following relationships between various quantities along the
Rankine-Hugoniot curve:
%
\begin{equation}
\begin{aligned}
  \frac{V_{02}}{V_{01}} &= \frac{(\Gamma+1)\,p_1 + (\Gamma-1)\,p_2}
                               {(\Gamma-1)\,p_1 + (\Gamma+1)\,p_2}
  \;\to\; \frac{\Gamma-1}{\Gamma+1}
  &&\text{for strong shock,} \\[6pt]
  %
  \frac{T_2}{T_1} &= \frac{p_2}{p_1}\,
    \frac{(\Gamma-1)\,p_2 + (\Gamma+1)\,p_1}
         {(\Gamma+1)\,p_2 + (\Gamma-1)\,p_1}
  \;\to\; \frac{\Gamma-1}{\Gamma+1}\,\frac{p_2}{p_1}
  &&\text{for strong shock,} \\[6pt]
  %
  \hat{\jmath}_0^{\,2} &=
    \frac{(\Gamma-1)\,p_1 + (\Gamma+1)\,p_2}{2V_{01}}
  \;\to\; \frac{\Gamma+1}{2}\,\frac{p_2}{V_{01}}
  &&\text{for strong shock,} \\[6pt]
  %
  f^2 &= \tfrac{1}{2}\,V_{01}
    \bigl[(\Gamma+1)\,p_1 + (\Gamma-1)\,p_2\bigr]
  \;\to\; \tfrac{1}{2}\,(\Gamma-1)\,V_{01}\,p_2
  &&\text{for strong shock,} \\[6pt]
  %
  \hat{\jmath}_0^{\,2} &=
    \frac{V_{01}\bigl[(\Gamma+1)\,p_1 + (\Gamma-1)\,p_2\bigr]^2}
         {2(\Gamma+1)}
  \;\to\; \frac{(\Gamma-1)^2}{2(\Gamma+1)}\,V_{01}\,p_2^{\,2}
  &&\text{for strong shock.}
\end{aligned}
  \tag{2.7.10}
\end{equation}
%
Here ``strong shock'' means ``in the limit $p_2/p_1 \gg 1$''.

\begin{figure}[htbp]
\centering
\fbox{\parbox{0.9\columnwidth}{\centering [Figure 2.7.2 placeholder]\\
The Poisson adiabat $\mathscr{P}\,\mathscr{P}'$ and the Rankine-Hugoniot
curve $\mathscr{R}\,\mathscr{R}'$ plotted in the $\mu V$--$p$ plane.}}
\caption{The Poisson adiabat $\mathscr{P}\,\mathscr{P}'$ and the
  Rankine-Hugoniot curve $\mathscr{R}\,\mathscr{R}'$ plotted in the
  $\mu V$--$p$ plane.}
\label{fig:2.7.2}
\end{figure}


%%%%%%%%%%%%%%%%%%%%%%%%%%%%%%%%%%%%%%%%%%%%%%%%%%%%%%%%%%%%%%%%
%%% SECTION 2.8  TURBULENCE
%%%%%%%%%%%%%%%%%%%%%%%%%%%%%%%%%%%%%%%%%%%%%%%%%%%%%%%%%%%%%%%%

\subsection{Turbulence}
\label{sec:2.8}

In the accretion of gas onto a black hole, turbulence probably plays an
important role.  (See \S\S4 and~5.)  Unfortunately, the theory of turbulence
is in a very uncertain state.  Little is known with confidence about the
astrophysical circumstances under which turbulence should develop, or about
the strength of the turbulence; see, e.g., Chapter~3 of
\citet{LandauLifshitz1959}; also \citet{Pikelner1961}.


%%%%%%%%%%%%%%%%%%%%%%%%%%%%%%%%%%%%%%%%%%%%%%%%%%%%%%%%%%%%%%%%
%%% SECTION 2.9  RECONNECTION OF MAGNETIC FIELD LINES
%%%%%%%%%%%%%%%%%%%%%%%%%%%%%%%%%%%%%%%%%%%%%%%%%%%%%%%%%%%%%%%%

\subsection{Reconnection of Magnetic Field Lines}
\label{sec:2.9}

\emph{Basic references}\quad \S5.3 of \citet{Cowling1965};
\citet{Sonnerup1970}; \citet{Yeh1970};
\citet{VainshteinZeldovich1972}.

\subsubsection*{Basic ideas}

In flat spacetime consider a plasma that is macroscopically neutral, that has a
high electrical conductivity~$\sigma$, that is sufficiently dilute for one to
ignore its dielectric properties: $\varepsilon = \mu = 1$.  Let the plasma be
endowed with a large-scale magnetic field, and examine the evolution of that
field in a Lorentz frame where the plasma has low velocity
$|\mathbf{v}| \ll c$.  Maxwell's equations for the electromagnetic field then
read
%
\begin{align}
  \nabla \times \mathbf{E} + (1/c)(\partial\mathbf{B}/\partial t) &= 0\,,
  \tag{2.9.1} \\[4pt]
  \nabla \times \mathbf{B} - (1/c)(\partial\mathbf{E}/\partial t)
    &= 4\pi\mathbf{J}/c\,.
  \tag{2.9.2}
\end{align}

In the local rest frame of the plasma the current is proportional to the
electric field, $\mathbf{J} = \sigma\mathbf{E}$.  When transformed to the
Lorentz frame where the plasma moves with velocity~$\mathbf{v}$, this
equation says
%
\begin{equation}
  \mathbf{J} = \sigma\bigl[\mathbf{E}
    + (\mathbf{v}/c) \times \mathbf{B}\bigr].
  \tag{2.9.3}
\end{equation}

A straightforward calculation from the above equations leads to the following
law for the rate of change of magnetic field along the world lines of the
plasma:
%
\begin{equation}
  \frac{d\mathbf{B}}{d\tau}
  = (\omega + \sigma - \tfrac{2}{3}\theta\,\mathbf{1}) \cdot \mathbf{B}
    - \frac{1}{4\pi\sigma}\left(
      \frac{\partial^2\mathbf{B}}{\partial t^2}
      - c^2\,\nabla^2\mathbf{B}\right).
  \tag{2.9.4}
\end{equation}
%
Here $\mathbf{1}$ is the unit 3-tensor; $\omega$, $\sigma$, and~$\theta$ are
the rotation, shear, and expansion of the plasma
[cf.\ eqs.~(2.5.17) and~(2.5.18)]; and
%
\[
  d/d\tau = \partial/\partial t + \mathbf{v}\cdot\nabla\,.
\]

Because of the very high electrical conductivity, one can ignore the
``wave-equation'' part of eq.~(2.9.4) almost everywhere in the plasma;
and the magnetic field evolves in a ``frozen-in'' manner:
%
\begin{equation}
  d\mathbf{B}/d\tau
  = (\omega + \sigma - \tfrac{2}{3}\theta\,\mathbf{1}) \cdot \mathbf{B}\,.
  \tag{2.9.5}
\end{equation}

[Cf.\ Eq.~(2.5.20) and associated discussion.]  However, in regions of plasma
where the field has strong gradients, the ``wave-equation'' part must come
into play.

Consider, as the case of greatest importance, a magnetic field that is chaotic
with an ordered structure on some scale~$l_c$ (subscript $c$ for ``cell
size'').  At the interface between adjacent cells, the magnetic field must
reverse sign (``neutral point'' or ``neutral sheet''), and its gradients will
be very large.  Hence, at the interface, ``wave-equation'' behavior can come
into play destroying the frozen-in behavior of the field.  The result is a
reconnection of magnetic field lines, as

\begin{figure}[htbp]
\centering
\fbox{\parbox{0.9\columnwidth}{\centering [Figure 2.9.1 placeholder]\\
(a)~Top view and (b)~perspective view of reconnecting magnetic field lines
in a plasma.  The reconnection region is stippled; the magnetic field lines
are drawn solidly and the flow lines of the plasma are dotted.  The innermost
field line is in the process of reconnecting into the dashed form.
The perspective view~(b) shows the closed curve~$\mathscr{E}$ around which
one integrates to derive eq.~(2.9.6) for the EMF in the connecting region.}}
\caption{The region in a plasma where magnetic field lines are
  reconnecting (schematic picture).  In the top view, (a), the
  reconnection region is stippled; the magnetic field lines are drawn
  solidly; and the flow lines of the plasma are dotted.  The innermost
  field line is in the process of reconnecting into the dashed form.
  The perspective view, (b), shows the closed curve~$\mathscr{E}$
  around which one integrates to derive eq.~(2.9.6) for the EMF in the
  connecting region.}
\label{fig:2.9.1}
\end{figure}
